\section{Conclusion}
\label{sec:conclusion}

In this paper, we introduced MSAE-AD, a novel few-shot anomaly detection framework that extends the minimalist SubspaceAD paradigm by replacing its linear PCA subspace with a Matryoshka Sparse Autoencoder (MSAE).
Our approach preserves the simplicity and efficiency of the original framework---frozen DINOv2 feature extraction followed by model fitting on normal samples---while introducing nonlinear, hierarchical sparse representations for more expressive normality modeling.

We evaluated two MSAE training strategies on the MVTec-AD and VisA benchmarks.
The reverse-weighted (RW) variant achieved average image-level AUROC scores of 98.72\% on MVTec-AD and 95.58\% on VisA, with pixel-level AUROC of 97.71\% and 98.46\%, respectively.
The uniform-weighted (UW) variant delivered competitive performance, achieving 98.07\% and 95.08\% image-level AUROC.
These results demonstrate that hierarchical sparse autoencoders offer a compelling alternative to linear subspace methods for few-shot anomaly detection.

Our per-category analysis revealed that MSAE-AD performs exceptionally well on structurally regular categories (achieving 100\% AUROC on several MVTec-AD classes) while facing challenges on categories with subtle, fine-grained anomalies.
The RW variant generally provides superior image-level detection, while UW shows advantages on specific categories with complex normal variation, suggesting that the optimal weighting strategy may be category-dependent.

Several directions for future work emerge from this study.
First, category-adaptive weighting strategies could combine the strengths of RW and UW by learning the optimal $\alpha$ coefficients from validation data.
Second, exploring higher expansion factors and deeper MSAE hierarchies (more granularity levels) may further improve representation capacity.
Third, the computational efficiency of MSAE training could be enhanced through optimized CUDA kernels enabling parallel multi-granularity processing.
Finally, investigating the complementary nature of PCA and MSAE anomaly scores---for example, through score-level fusion---could yield a hybrid approach that leverages the strengths of both linear and nonlinear modeling.

In summary, MSAE-AD provides a principled and effective extension of the SubspaceAD framework, demonstrating that hierarchical sparse representations can serve as powerful normality models for industrial anomaly detection.
We hope this work encourages further exploration of representation learning techniques for practical, data-efficient visual inspection systems.
